\section{Introduction}
\label{sec:introduction}

Realized-volatility forecasting sits in an unusual corner of the prediction landscape. The target is among the most forecastable objects in finance---simple heterogeneous-autoregressive (HAR) models on lagged realized variance explain the large majority of out-of-sample variation \citep{Corsi2009, AndersenBollerslevDieboldLabys2003}---yet the increments of skill beyond that autoregressive core are small, fragile, and easy to manufacture spuriously. Machine-learning comparisons in this literature routinely disagree about whether nonlinear models help at all \citep{christensen2023machine, gu2020empirical}, and we will show that part of that disagreement is mechanical: rankings between model classes flip with hyperparameter defaults, evaluation panels, and imputation conventions long before they reflect anything about volatility.

This paper is an attempt to make those increments of skill legible. We work on a single fixed panel---218{,}934 thirty-minute bars of S\&P~500 realized variance spanning 2005--2024, with 41 exogenous predictors organized into eight economic buckets---and we hold one benchmark fixed throughout: an ordinary-least-squares HAR model with calendar features, refit every bar, with out-of-sample QLIKE of $0.13415$. Every model in the paper is evaluated on the identical row index, so every comparison is bar-for-bar, and every claim of improvement carries a Diebold--Mariano statistic with heteroskedasticity-and-autocorrelation-consistent inference \citep{DieboldMariano1995, HarveyLeybourneNewbold1997} or a model-confidence-set membership \citep{HansenLundeNason2011}.

The organizing empirical fact, developed in Section~\ref{sec:descriptive}, is that the exogenous signal is \emph{dense but weak}. The design matrix has 529 columns whose covariance concentrates in a handful of directions, but the forecasting signal does not live in those directions: the first principal component carries essentially none of it, low-rank compressions destroy it, and $\ell_1$ selection recovers it only as a small unstable core plus a wide shallow tail. No individual feature is strong; many features are weakly and redundantly informative. Each subsequent section asks what this regime implies for a different modeling decision.

Section~\ref{sec:marginal} measures where the information lives, using only the least-squares model class so that attribution is not confounded with model capacity. Every bucket improves on HAR and every improvement is statistically decisive, but the gains overlap heavily: the sum of individual bucket contributions is roughly twice the joint contribution, and the model confidence set over feature sets is the singleton containing the joint bucket. A Frisch--Waugh--Lovell block decomposition attributes 96\% of explainable variance to the six HAR features, and an exact $2^8$ factorial over individual feature bundles finds one dominant cheap feature, one partially redundant companion, and a long tail of tiny effects.

Section~\ref{sec:linvsnonlin} asks whether extracting the signal requires nonlinearity. On a battery of five estimators crossed with nine feature sets---all causally tuned, all on the shared panel---tree ensembles hold a modest but systematic edge over penalized linear models, and the edge decomposes cleanly: a pure-nonlinearity premium of $0.00297$ visible on HAR features alone, on top of the information premium. Crucially, the edge is not a hyperparameter story. Causally tuned hyperparameters trail fixed defaults; an oracle over penalty configurations improves the incumbent by at most $0.00015$; and the nonlinear capacity that matters saturates at additive-plus-pairwise structure anchored to the session clock.

Section~\ref{sec:algorithms} then freezes the feature set entirely and treats the estimation \emph{algorithm} as the design space. We present four designs: exact rank-1 rolling least squares (the incumbent's own machinery), a recursive elastic net via online homotopy---the $\ell_1$ analog of the Sherman--Morrison update \citep{Garrigues2008}---a causal rolling hyperparameter-selection protocol, and, as the centerpiece, a spectral $k$-nearest-neighbors estimator that embeds trailing residual trajectories with Laplacian eigenmaps \citep{BelkinNiyogi2003} and forecasts by analog retrieval. The section's honesty is its contribution: ablations show that the manifold-embedding step hurts on every matched configuration, while the OLS-anchored short-window analog search that remains beats the incumbent decisively.

Three methodological commitments run through the paper. First, one panel, one benchmark: no result is quoted across incompatible evaluation regimes, and each table names its panel. Second, significance before narrative: differences without a DM statistic or MCS membership are labeled as unresolved. Third, negative results are reported at the same altitude as positive ones---the failures of PCA compression, deep sequence models, hyperparameter tuning, and the manifold embedding are load-bearing facts about the dense-but-weak regime, not footnotes.
